\scanpage{77}
précisément pour que la \deleted{\unsure{factorisation}} \inserted{réalisation} géom
obtenue de $\Pi_n$ $(3\leq n\leq\infty)$ soit
fidèle. De cette façon, on trouvera en
particulier la classification des \deleted{\unsure{po|formes}} $n$-polygônes
réguliers sur $k$, indépendamment de la
donnée d'un $\Pi$-épinglage, par des
\deleted{\unsure{élé}} un invariant $\alpha\in k$ satisfaisant aux \unsure{mêmes}
conditions.

c) Déduire de ceci les \deleted{des} propriétés
\unsure{essentielles|matricielles} \deleted{des} (en particulier la classification)
des $n$-polygônes réguliers, quel que soit $n$
et indépendamment de tt épinglage
par un repère. (Cette étape est
essentiellement triviale \ldots)

\uline{2}. Reformulation du Pb des réalisations géom
\begin{tabular}{|p{.91\linewidth}}régulières de $\Pi_\infty$ - dans plan affine $E$
(sur \uline{corps} $k$, pour simplifier \ldots)\end{tabular}

Rappelons $G=\operatorname{Aut}(\Pi_\infty)\simeq D_\infty\inserted{{}=\unsure{Ga_1|G_{a1}|G_1|G_{0,1}}}\simeq\{\sigma_0,\sigma_1\mid\sigma_0^{\unsure{2|7|1}}=\deleted{\unsure{\sigma}}\sigma_1^2=1\}$

a) Représentation
\[ \varphi_G:D_\infty\longrightarrow\operatorname{Aff}(E) \tag*{(1)} \]
% archive sheet number: 76
